\section{What a multiplicity bound would prove}
\label{sec:parameters}

We give a rigorous parameter bound for one fixed-order scan that fully
cancels after every crossing.  It applies to the ordinary and pointed
engines; deferred tails and races are excluded from this statement.  The
input is a normalized $n$-crossing PD code, so edge labels use $O(\log(n+1))$
bits.  The crossing-free case is handled separately in constant time.

Let $w$ be the maximum number of frontier endpoints, including the final
two endpoints of a pointed scan.  Let $M\geq1$ be the maximum number of live
objects \emph{after} a fully canceled stage, including the initial stage.
Let $\mu$ be the largest total multiplicity of any one matching at any such
stage, summed over all homological degrees.  These are parameters of the
actual computation; none is assumed to have a small bound in terms of $n$.

\begin{lemma}[One-crossing expansion]
\label{lem:eightfold}
If a canceled stage has at most $M$ objects, attaching the next crossing and
delooping creates at most $8M$ objects before cancellation.  At most $4M$
cancellations can occur at that stage.
\end{lemma}

\begin{proof}
There are two smoothings.  Every newly closed circle contains at least one
of that smoothing's two new arcs: previously closed circles have already
been delooped, and untouched old matching arcs cannot form a new circle.
Thus each smoothing creates at most two closed circles.  Delooping supplies
at most $2^2$ copies, giving at most $2\cdot2^2=8$ new objects per old object.
Taking the pointed quotient cannot increase this count.  Every cancellation
removes two objects and creates none, which gives the second assertion.
\end{proof}

This argument includes the final crossing, whether it closes the ordinary
diagram completely or leaves the two marked endpoints.  It does not apply
repeatedly with the same $M$ through an unreduced tail: $r$ crossings without
intermediate cancellation permit the weaker bound $8^rM$.  The parameter
must then be enlarged accordingly.

\begin{proposition}[Work in terms of retained objects]
\label{prop:retainedbound}
The full-elimination scan performs at most $O(nM^3)$ algebra compositions.
Its lazy pivot queue uses at most $O(nM^4)$ insertions and removals.  With
deterministic dictionary implementations, its bit-operation complexity is
bounded by
\begin{equation}
 \poly(n,w,\log(M+1))\,M^4\,2^{O(w)}.
 \label{eq:retainedbound}
\end{equation}
For the supplied Python dictionaries, the same expression is an operation
bound in the usual dictionary cost model; constant worst-case hashing cost
is not being asserted.
\end{proposition}

\begin{proof}
Put $N=8M$.  A stage has at most $N^2$ differential entries.  Transferring an
entry through either smoothing produces at most sixteen delooped entries,
so crossing attachment involves $O(M^2)$ bounded-arity transfer operations.
Each pivot has at most $N$ incoming and $N$ outgoing neighbors.  Its Schur
update therefore uses $O(N^2)$ compositions.  There are at most $N/2$ pivots,
giving $O(M^3)$ compositions per stage.

The implementation's lazy queue needs separate accounting.  Initialization
and pivot updates create $O(M^3)$ candidate records per stage, with duplicates
allowed.  A stale record is refiled at its newly evaluated Markowitz cost.
It cannot become stale again until a cancellation changes the differential.
Since there are $O(M)$ cancellations, each created record is refiled at most
$O(M)$ times.  Thus there are $O(M^4)$ queue operations, each with at most
logarithmic heap overhead.  This argument avoids incorrectly equating cubic
composition count with cubic total bookkeeping cost.  An indexed queue or
complete pivot search can achieve a different accounting, but is not the
current implementation.

An overlay on at most $w$ endpoints has at most $w/2$ circles.  An unpointed
morphism consequently has at most $2^{w/2}$ coefficient bits; the pointed
quotient removes one variable when active.  Enumerating pairs of monomials,
compiling their surface components, and evaluating each component all take
$2^{O(w)}\poly(w)$ bit operations, even with naive integer operations.
Labels, object indices, and matching identifiers add logarithmic factors in
$n$ and $M$.  Numerical cache entries are bounded by the actual polynomial
number of calls just counted, rather than by a count of possible values:
a Hom space with $2^{w/2}$ coefficient bits has up to
$2^{2^{w/2}}$ possible coefficient vectors.  Finitely many topology-plan
shapes therefore do not remove the multiplicity parameter.  Replacing
hashing by balanced dictionaries adds only logarithmic factors in the
number of stored entries.  Combining these observations over $n$ stages proves
\eqref{eq:retainedbound}.
\end{proof}

\begin{corollary}[Width and multiplicity criterion]
\label{cor:multiplicitycriterion}
Define $D(0)=1$ and $D(w)=(w-1)!!$ for even $w\geq2$.  Then
\[
 M\leq\mu D(w),\qquad
 T\leq\poly(n)\exp\!\bigl(O(w\log(w+1)+\log(\mu+1))\bigr).
\]
In particular, a family satisfying
$w\log(w+1)+\log(\mu+1)=O(\log^2(n+2))$
has a quasi-polynomial implementation of this scan.
\end{corollary}

\begin{proof}
On a fixed frontier of $b$ endpoints there are at most $(b-1)!!$ perfect
matchings.  Multiplying by the total multiplicity bound proves the first
inequality.  This is a \emph{per-stage} count: different stages have different
label sets, and retaining their matchings adds at most a factor $n+1$.
Finally $\log D(w)=O(w\log(w+1))$, and substitution into
\eqref{eq:retainedbound} gives the claimed form.
\end{proof}

The double factorial is deliberate.  A Catalan count requires additional
common-boundary topology, discussed earlier; arbitrary frontier sets do not
supply it.  If every frontier of the \emph{actual scan} has a certified common
disk realization, including the cut endpoints when a pointed engine is used,
the sharper count gives
\[
 M\leq\mu C_{w/2},\qquad
 T\leq\poly(n)\exp\!\bigl(O(w+\log(\mu+1))\bigr),
\]
because the Catalan number satisfies $C_{w/2}\leq2^w$.
Then $w+\log(\mu+1)=O(\log^2(n+2))$ suffices for quasi-polynomial time.
The last-crossing marking theorem preserves frontier cardinality; it does
not by itself certify this common disk hypothesis.

More seriously, small width does not force small multiplicity.
Serial diagrams of connected sums of $k$ trefoils have bounded scan width
and reduced rank $3^k$.  The trefoil's characteristic-two reduced complex
has three surviving generators; the connected-sum tensor product multiplies
these ranks~\cite[Sections~3--4]{khovanov2002}.  Processing one fixed-size factor at a time keeps
only a bounded number of frontier endpoints.  Nevertheless the final
pointed minimal complex has $3^k$ objects of the same matching.  This is an
exponential obstruction for the \emph{explicit generator representation}.
It is not a lower bound for computing the integer $3^k$, a compressed graded
answer, or the unknot decision.  Factorization and the isolated-summand
criterion can avoid that explicit representation on this family.

\begin{corollary}[Small visible factors]
\label{cor:small-factors}
Suppose the atomic visible factors of a validated $n$-crossing diagram have
at most $b$ crossings each.  Checked factorization followed by complete
scanning of each factor gives an exact decision procedure with bound
\[
 O(n\log(n+1))+n\,2^{O(b)}.
\]
In particular, $b=O(\log^2(n+2))$ gives the sharper quasi-polynomial target.
\end{corollary}

\begin{proof}
Lemma~\ref{lem:eightfold} bounds all object counts on an $s$-crossing factor
by $8^s$, even if no cancellation helps.  Its frontier has $O(s)$ endpoints.
Proposition~\ref{prop:retainedbound} therefore gives $2^{O(s)}$ time.
There are at most $n$ nonempty factors, their crossing counts sum to $n$,
and a connected sum is the unknot precisely when all factors are.
Add the checked factorization cost from
Proposition~\ref{prop:interlace-reconstruction}.
\end{proof}

This restricted class already has a quasi-polynomial algorithm with a
polynomial-cost factorizer.  The present improvement reduces its
preprocessing cost; it does not by itself enlarge the class to all diagrams.

\subsection{A further closed-complex rank obstruction}

The following elementary observation is a proved, unimplemented follow-up,
without a novelty-priority claim.  It can sometimes avoid final scalar
elimination when no isolated object is present.

\begin{lemma}[Vertex-cover bound]
\label{lem:vertexcoverbound}
Let $(V,d)$ be a finite scalar chain complex with $N$ homogeneous basis vectors.
Form an undirected graph joining two basis vectors whenever the differential
has a nonzero entry between them.  For every vertex cover $C$,
\[
 \dim H(V,d)\geq N-2|C|.
\]
\end{lemma}

\begin{proof}
The complementary vertices $I$ are independent, so
$d(\operatorname{span}I)\subseteq\operatorname{span}C$.
Put $b=\dim d(\operatorname{span}I)$.  Since $d^2=0$, this $b$-dimensional
subspace lies in the kernel of $d$ restricted to $\operatorname{span}C$.
Thus $\dim d(\operatorname{span}C)\leq|C|-b$, and
$\operatorname{rank}d\leq b+(|C|-b)=|C|$.
Now use $\dim H=N-2\operatorname{rank}d$.
\end{proof}

The support graph is bipartite by homological parity, so a minimum cover
can be found by bipartite matching.  A cover giving $N-2|C|>1$ proves
knottedness for a closed reduced knot complex.  Verifying the cover against
that scalar differential is easy, although obtaining and verifying the
differential still requires its preceding computation.  The square-zero
hypothesis is essential, and the lemma cannot be applied to an open
cobordism complex simply by counting its matching objects.
